\chapter{D2NN Beam-Reducer Sweep}

\section{Motivation}

The first successful track in the chronology was a spatial D2NN beam-reducer
receiver. The aim was pragmatic: if the telescope output beam can be compressed
into a detector-friendly footprint while preserving most of the energy, then a
compact phase-only network might already provide value before any Fourier-domain
architecture is considered.

\section{Experimental Setup}

The run family \texttt{0325-telescope-sweep-cn2-5e14-15cm} compares five loss
formulations with a common five-layer spatial D2NN. All variants use the same
geometry and differ only in training loss. The committed summary shows that the
best complex-overlap results cluster tightly near a plateau around $0.646$,
with throughput effectively unity in every case.

\begin{table}[H]
  \centering
  \begin{tabular}{lcccc}
    \toprule
    \smalltablehead{Variant} & \smalltablehead{$\CO$} & \smalltablehead{IO} & \smalltablehead{EE} & \smalltablehead{$\TP$} \\
    \midrule
    baseline\_co & 0.6456 & 0.7843 & 0.8399 & 1.0000 \\
    co\_amp & 0.6456 & 0.7844 & 0.8400 & 1.0000 \\
    co\_ffp & 0.6407 & 0.7149 & 0.7709 & 1.0000 \\
    co\_phasor & 0.6295 & 0.7647 & 0.8186 & 1.0000 \\
    roi80 & 0.5457 & 0.8074 & 0.9880 & 1.0000 \\
    \bottomrule
  \end{tabular}
  \caption{\texttt{0325} beam-reducer sweep summary. The best field-level score
  saturates around $\CO \approx 0.646$, while detector-like concentration
  metrics can be traded against field fidelity.}
\end{table}

\runfigure{0325-telescope-sweep-cn2-5e14-15cm/phase2_fig1_input.png}
  {0325 input-field summary}
  {Representative input-field visualization for the 15\,cm telescope receiver
  experiment family.}

\runfigure{0325-telescope-sweep-cn2-5e14-15cm/phase2_fig2_d2nn_output.png}
  {0325 output-field summary}
  {D2NN output-field comparison across the main \texttt{0325} loss variants.}

\runfigure{0325-telescope-sweep-cn2-5e14-15cm/phase2_fig3_detector.png}
  {0325 detector summary}
  {Detector-plane energy redistribution in the beam-reducer sweep.}

\runfigure{0325-telescope-sweep-cn2-5e14-15cm/phase2_fig4_masks.png}
  {0325 masks}
  {Learned phase masks for the five-layer spatial D2NN beam reducer.}

\runfigure{0325-telescope-sweep-cn2-5e14-15cm/phase2_fig5_training_curves.png}
  {0325 training curves}
  {Training trajectories for the beam-reducer sweep.}

\runfigure{0325-telescope-sweep-cn2-5e14-15cm/phase2_fig6_statistics.png}
  {0325 statistics}
  {Aggregate statistics showing the plateau in field-level performance.}

\section{Interpretation}

Two lessons emerged immediately. First, the spatial D2NN could shape the
detector-plane distribution without catastrophic throughput loss. Second, the
best field-level score plateaued early. In other words, the beam reducer was
useful but not obviously approaching a vacuum-field recovery regime.

\begin{discovery}
The \texttt{0325} beam-reducer receiver demonstrated that passive diffractive optics
can improve detector-side concentration while preserving energy, but it also
showed an empirical plateau in field-level quality. That plateau motivated the
shift toward FD2NN architectures and later detector-centric objective design.
\end{discovery}

The next step was therefore not incremental tuning of the same spatial design.
Instead, the campaign moved toward a Fourier-domain architecture in the hope
that a 4f representation might offer stronger control over wavefront structure.
